\documentclass{article}
\usepackage{graphicx} % Required for inserting images
\usepackage[english]{babel}

\usepackage[letterpaper,top=2cm,bottom=2cm,left=3cm,right=3cm,marginparwidth=1.75cm]{geometry}

\usepackage{amsmath}
\usepackage{graphicx}
\usepackage[colorlinks=true, allcolors=blue]{hyperref}

\title{Data 501 Literature Review II}
\author{Ashton Frese, Karen Reyes González, Wilbur Elbouni}
\date{February 24, 2026}

\begin{document}

\maketitle 
\title{Paper Review: \textit{The Transmission of Shocks Among S\&P Indexes}}

\section{Paper Overview}


\subsection{Citation:}  
B. T. Ewing, ``The Transmission of Shocks Among S\&P Indexes,'' \textit{Applied Financial Economics}, vol. 12, pp. 285--290, 2002.

\subsection{Problem Addressed:}  
The paper investigates the interrelationship between the sector indexes of the U.S. equity market. It particularly focuses on the relation between the five S\&P sector stock price indexes and the proportion of forecast error variance for each index that can be explained by the innovations in the other indexes.


\subsection{Motivation:}  
With the increasing popularity of index investing and the utilization of the market and sector indexes as a basis for mutual fund and institutional portfolios, it has become important for investors to know:

\begin{itemize}
    \item how strongly sector indexes co-move,
    \item which sectors are more influential in driving movements in others,
    \item and what this implies for diversification across sector-based instruments.
\end{itemize}

\noindent The author argues that understanding the transmission of shocks across sector indexes informs both index fund design and sector allocation decisions, since cross-index linkages constrain the benefits of diversification.

\hfill

\section{Methodology Summary}

\subsection{Data:}  
The study uses monthly U.S. data over the post-1987 crash period, from January 1988 to July 1997, on five Standard \& Poor's sector stock price indices:

\begin{itemize}
    \item Capital Goods (LNCAP),
    \item Financials (LNFI),
    \item Industrials (LNIN),
    \item Transportation (LNTR),
    \item Utilities (LNUT).
\end{itemize}

\noindent All index series are converted to natural logarithms and then first-differenced to obtain log-returns (denoted by $\Delta$), based on unit root tests that indicate non-stationarity in levels but stationarity in first differences.

\subsection{Techniques and Models:}  
The paper estimates a five-variable Vector Autoregression (VAR) of order one:

\[
\Delta r_t = A_0 x_t + u_t,
\]

\noindent where $\Delta r_t$ is the vector of  one of the five sector log-returns $(\Delta \text{LNUT}, \Delta \text{LNTR}, \Delta \text{LNIN}, \Delta \text{LNFI}, \Delta \text{LNCAP})^\top$, and $x_t$ stacks the lagged returns. Lag length is chosen via Akaike Information Criterion (AIC), Schwarz Bayesian Criterion (SBC), and likelihood ratio tests.

\hfill

\noindent A key methodological feature is the use of \textbf{generalized forecast error variance decomposition} (GFVD), following Koop et al.\ (1996) and Pesaran and Shin (1998), rather than traditional orthogonalized decompositions based on a Cholesky factorization. The generalized approach:

\begin{itemize}
    \item allows for contemporaneous correlation among shocks,
    \item is invariant to the ordering of variables in the VAR,
    \item and provides meaningful decompositions even at short horizons (including initial impact).
\end{itemize}

\noindent More formally, the VAR is specified in its moving average representation, and for each $i$th index and $j$th source of shock, the generalized decomposition $\Theta_{ij,N}$ measures the fraction of the $N$-step ahead forecast error variance of $i$ due to shocks in $j$, relative to the total forecast error variance in $i$ at horizon $N$.

\hfill

\noindent Prior to the application of GFVD, the author tests for contemporaneous correlation of shocks across equations using a log-likelihood ratio test of the covariance matrix of the residuals. The null hypothesis that off-diagonal elements of the matrix are zero is strongly rejected.

\hfill

\section{Key Contributions}

\subsection{Main Findings:}  
The generalized forecast error variance decompositions are presented for forecast horizons up to three months, after which they settle. The main findings are as follows:

\begin{itemize}
    \item \textbf{Utilities (Panel A):}  
    The forecast variance of the utilities index ($\Delta$ LNUT) is driven by shocks in other industries as follows:    
    \begin{itemize}
        \item Financials explain around 26--28\% of the forecast variance.
        \item Industrials explain around 20--21\%.
        \item Capital Goods explain approximately 13--17\%.
        \item Transportation explains about 10--12\%.
    \end{itemize}
    Utilities are thus strongly influenced by the rest of the market.

    \item \textbf{Transportation (Panel B):}  
    Transportation returns ($\Delta$LNTR) are strongly influenced by Capital Goods and Industrials:
    \begin{itemize}
        \item Capital Goods explains roughly 58--64\% of forecast variance.
        \item Industrials explains around 52--57\%.
        \item Financials explain about 37--38\%.
        \item Utilities explain around 10--13\%.
    \end{itemize}

    \item \textbf{Industrials (Panel C):}  
    For Industrials ($\Delta$LNIN), Capital Goods dominate:
    \begin{itemize}
        \item Capital Goods explains about 84--86\% of forecast variance.
        \item Financials and Transportation each explain roughly 57\%.
        \item Utilities contribute around 21--22\%.
    \end{itemize}

    \item \textbf{Financials (Panel D):}  
    Financials ($\Delta$LNFI) are explained mainly by Industrials and Capital Goods:
    \begin{itemize}
        \item Industrials explain about 53--57\%.
        \item Capital Goods explain around 49--52\%.
        \item Transportation explains roughly 37\%.
        \item Utilities contribute about 28--30\%.
    \end{itemize}

    \item \textbf{Capital Goods (Panel E):}  
    For Capital Goods ($\Delta$LNCAP):
    \begin{itemize}
        \item Industrials explain about 85--86\% of forecast variance.
        \item Transportation explains around 64--65\%.
        \item Financials explain roughly 51--52\%.
        \item Utilities contribute around 13\%.
    \end{itemize}
\end{itemize}

\noindent Across all panels, Utilities are the least influential in explaining the forecast error variance of other indexes, while Capital Goods and Industrials are found to be the most dominant in shock transmission to the rest of the system.

\subsection{Innovations and Insights:}

\begin{itemize}
    \item The paper is among the first in financial economics to apply generalized forecast error variance decompositions, avoiding the arbitrary ordering problem inherent in orthogonalized (Cholesky-based) approaches.
    \item It provides a clear ranking of sector indexes in terms of their \textit{explanatory strength}: sectors like Capital Goods and Industrials serve as ``shock transmitters'', while Utilities appear more ``downstream''.
    \item The decompositions quantify the limits of diversification across S\&P sector indexes, since a substantial portion of each sector’s variance is explained by shocks originating in other sectors.
\end{itemize}

\hfill

\section{Critical Reflection}

\subsection{Strengths:}

\begin{itemize}
    \item \textbf{Sector-Level Interdependence:}  
    Examining intra-country, cross-sector linkages (as opposed to international linkages or a composite index) enables the paper to tackle issues related to sectoral diversification and contagion in the US sectors directly.
    
    \item \textbf{Robust Time-Series Decomposition:}  
    The motivation behind generalized variance decompositions is well-placed in the context of strong contemporaneous correlations in the VAR residuals. In addition, the ordering-invariance property is especially relevant at short horizons. 
    
    \item \textbf{Interpretability:}  
    The decomposition methodology provides a very intuitive measure: "percentage of forecast variance of sector A explained by shocks in sector B," which is easily interpretable in terms of portfolio construction.
\end{itemize}

\hfill

\subsection{Limitations and Assumptions:}

\begin{itemize}
    \item \textbf{No Explicit Macroeconomic Drivers:}  
    Unlike Ewing et al. (2003), this study does not include macroeconomic factors like output, monetary policy, or risk premia. The only factors included in this study are those from the sector index system. Therefore, there is no explanation for the occurrence of these factors.
    
    \item \textbf{Data Window and Frequency:}  
    The data used in this study only covers the period from 1988 to 1997. The sample frequency used is only at a monthly level. Therefore, it excludes major crises occurring after 1997, like the 2008 crisis, the COVID-19 crisis, and the 2022 tightening crisis. It also excludes higher frequency data, like daily data, which might be useful for short-run risk management and event study.
    
    \item \textbf{No Forecasting Performance Assessment:}  
    The VAR model used in this study is only descriptive. It decomposes the variance of the sector index. It is not used for forecasting purposes. Therefore, it is not compared with other models using metrics like Root Mean Squared Error or Mean Absolute Error.
    
    \item \textbf{Static Reporting, No Visual Analytics:}  
    The results obtained from the VAR model are only reported in a single table. There is no exploration of the results using visual analytics.
\end{itemize}

\hfill

\section{Connection to Our Project}

\noindent This paper is complementary to our capstone project as it examines linkages between sectors in the U.S. equity market, which is highly relevant to our work on sector ETFs and rolling correlations.

\subsection{Conceptual Alignment:}  
Our project seeks to develop an interactive dashboard for analyzing and forecasting sector performance and volatility, with capabilities to:

\begin{itemize}
    \item compare sectors (Technology, Energy, Financials, Healthcare, Industrials), 
    \item understand how shocks and events propagate across sectors,
    \item evaluate the diversification potential during different regimes (shock vs.\ stable periods).
\end{itemize}

\noindent Ewing (2002) breaks down the explanation of variance in other S\&P sector indexes using VAR, which supports our decision to model and visualize co-movement at the sector level instead of modeling sectors as separate series.


\subsection{Influence on Our Analytical Design}

\begin{itemize}
    \item \textbf{Motivating Rolling Correlations and Network Views:}  
    The paper’s GFVD results effectively define a static ``influence network'' among sectors (e.g., Capital Goods and Industrials as central transmitters, Utilities as peripheral). In our dashboard, we plan to:
    \begin{itemize}
        \item compute rolling pairwise correlations or dependence measures among sector ETF returns.
        \item potentially visualize sector linkages as a dynamic network or heatmap.
    \end{itemize}
    \noindent Ewing’s decompositions provide a conceptual baseline for what such a network should reveal in terms of central vs.\ peripheral sectors.

    \item \textbf{Internal vs.\ External Shocks:}  
    Whereas our other reviewed paper (Ewing et al., 2003) focuses on macroeconomic shocks (output, monetary policy, risk premia), this paper focuses on \textit{internal} sector-to-sector transmission. Together, they motivate a two-layer view in our project:
    \begin{enumerate}
        \item macro events and policy changes as external drivers,
        \item and cross-sector linkages as the internal propagation mechanism visible in ETF returns.
    \end{enumerate}
\end{itemize}


\subsection{Methodological Inspiration and Differences:}  

\begin{itemize}
    \item \textbf{From GFVD to Exploratory Diagnostics:}  
    Although we do not necessarily carry out a full GFVD, the concept of decomposing forecast error variance into the contribution of sector i, can be approximated through:
    \begin{itemize}
        \item Rolling regression or VAR approaches,
        \item Exploratory methods such as rolling R2 statistics from regressions of sector i on other sectors.
    \end{itemize}
    These exploratory diagnostics can be incorporated into the dashboard to allow users to investigate which sectors are most strongly connected.

    \item \textbf{Integration with Forecasting Models:}  
    Ewing (2002) applies a VAR solely for decomposition purposes. For our project, we will concentrate on integrating forecasting models (ARIMA, GARCH, LSTM) with sector ETFs and volatility. We could still apply VAR-style analysis to:
    \begin{itemize}
        \item determine whether to assume sectors are conditionally independent in our forecasting framework,
        \item or whether to extend to multi-variate models (such as VAR-GARCH or multivariate LSTM) in future work, if there is evidence of strong cross-sector linkages
    \end{itemize}
\end{itemize}

\subsection{Extension of Scope:}  
We build upon Ewing (2002) in several key ways:

\begin{itemize}
    \item \textbf{Data Frequency and Horizon:} From monthly S\&P sector indices (1988-1997) to daily U.S. sector ETFs (2005-2026), including recent crises.
    \item \textbf{External Drivers:} From index-based data shocks to a hybrid model using sector ETFs, Fed Funds Rate, and curated macro/geopolitical events.
    \item \textbf{Models and Evaluation:} From descriptive VAR approaches \& decomposition analyses to a predictive framework (ARIMA, GARCH, LSTM) \& evaluation through formal metrics (RMSE/MAE) across regimes of shock \& stability.
    \item \textbf{Interface:} From static decompositions to interactive visualizations using libraries like Plotly/Dash for exploring sector transmission patterns for non-expert users.
\end{itemize}

\noindent In conclusion, Ewing (2002) provides a strong theoretical framework for our capstone project in several ways. First, it reinforces the significance of sector indexes through a demonstration of the interconnection between them. This supports our use of sector ETFs in our project. Second, it highlights the structured nature of sector shocks and their propagation, which is directly connected to our project.

\end{document}